\documentclass[11pt]{article}

\usepackage[T1]{fontenc}
\usepackage[spanish,es-nodecimaldot]{babel}
\usepackage{amsmath,amssymb,mathtools,bm}
\usepackage{booktabs}
\usepackage{geometry}
\usepackage{enumitem}
\usepackage[hidelinks]{hyperref}

\geometry{margin=1in}

\newcommand{\R}{\mathbb{R}}
\newcommand{\D}{\Delta}
\newcommand{\argmin}{\operatorname*{arg\,min}}

\title{Desarrollo matematico del seguimiento de minimos locales\\
para seleccionar suavidad y orden de tendencia}
\author{Nota tecnica del proyecto \texttt{paper\_numerical-smoothness-selection}}
\date{}

\begin{document}
\maketitle

\begin{abstract}
Esta nota formaliza, paso a paso, el procedimiento exploratorio implementado en
\texttt{experiments/numerical\_smoothness\_selection/run\_two\_stage\_order\_validation.py}.
La idea central es separar dos tareas. Primero, una perdida de
\emph{Validation 1} define la geometria de candidatos: para cada orden de
diferencia \(d\), se identifican varios minimos locales de la funcion de
perdida en la coordenada de suavidad normalizada \(S\in[0,1]\). Segundo, esos
minimos se siguen a traves de ventanas temporales cercanas y cada rama se
evalua en un bloque posterior de \emph{Validation 2}. La rama final se escoge
por persistencia temporal y desempeno historico en Validation 2, no por la
misma perdida que la genero. Finalmente, la rama seleccionada se actualiza una
vez mas en una Validation 1 inmediatamente anterior al test, y el test
verdadero permanece sin usar hasta el final.

El punto importante es que el valor final de \(S\) \textbf{no} es un promedio
global de suavidades. Es el minimo local mas reciente perteneciente a una rama
que ha sido seguida localmente en el tiempo y cuya calidad predictiva fue
evaluada fuera de la ventana que genero cada minimo.
\end{abstract}

\tableofcontents

\section*{Nota numerica complementaria}
La derivacion detallada de las primeras y segundas derivadas, la inversion
\(S\leftrightarrow\lambda\), el refinamiento adaptativo de intervalos, la
construccion de brackets, el metodo de Brent, la clasificacion por curvatura y
el tratamiento de endpoints se desarrolla en
\texttt{numerical\_methods\_smoothness\_search.tex}. Esta nota principal se
concentra en el tracking temporal y en la seleccion final de la rama.

\section{Objetivo y notacion}

Sea una serie positiva
\[
Y_1,\ldots,Y_N,\qquad Y_t>0,
\]
y definamos
\[
X_t=\log Y_t.
\]
En los experimentos aplicados se trabaja sobre \(X_t\). Cuando la metrica de
seleccion es \texttt{level\_rmse}, los pronosticos se transforman de regreso a
nivel mediante la exponencial antes de calcular el error.

Los parametros estructurales son:
\begin{itemize}[leftmargin=2em]
    \item \(d\in\{1,2,3,4\}\): orden de la penalizacion por diferencias;
    \item \(L\): longitud de la ventana de entrenamiento;
    \item \(h\): horizonte de Validation 1, Validation 2 y test final;
    \item \(S\in[0,1]\): coordenada normalizada de suavidad;
    \item \(\lambda\in[0,\infty]\): parametro clasico de penalizacion;
    \item \(\varepsilon_{\mathrm{track}}>0\): radio local maximo para continuar
    un minimo de una ventana a la siguiente.
\end{itemize}

En el codigo actual,
\[
\varepsilon_{\mathrm{track}}=0.10,
\]
se inicializan como maximo cinco ramas por orden \(d\), y se usa una separacion
numerica de \(0.02\) para evitar contar dos veces el mismo minimo en una misma
superficie. Para GDP se usa \(h=8\) trimestres; para SPY, AAPL y BTC--USD se
usa \(h=60\) observaciones diarias. El avance de los rolling origins es de una
observacion para GDP y de cinco observaciones para las series diarias.

\section{Suavizador penalizado}

\subsection{Problema de minimos cuadrados penalizados}

En una ventana de entrenamiento de longitud \(L\), sea
\[
\bm{x}=(x_1,\ldots,x_L)^\top\in\R^L.
\]
Sea \(D_d\) la matriz de diferencias finitas de orden \(d\). La tendencia
penalizada es
\begin{equation}
\widehat{\bm{\tau}}_{\lambda}
=
\argmin_{\bm{\tau}\in\R^L}
\left\{
\|\bm{x}-\bm{\tau}\|_2^2
+
\lambda\|D_d\bm{\tau}\|_2^2
\right\}.
\label{eq:pls}
\end{equation}
Definiendo \(Q_d=D_d^\top D_d\), la solucion es
\begin{equation}
\widehat{\bm{\tau}}_{\lambda}
=
H_{\lambda,d}\bm{x},
\qquad
H_{\lambda,d}=(I+\lambda Q_d)^{-1}.
\label{eq:smoother}
\end{equation}

Cuando \(\lambda=0\), \(H_{0,d}=I\). Cuando \(\lambda\to\infty\), sobreviven
unicamente las componentes del null-space:
\[
H_{\lambda,d}\longrightarrow P_{\ker(D_d)}.
\]
Por tanto, la tendencia limite tiene diferencia \(d\)-esima nula.

\subsection{Coordenada normalizada de suavidad}

Sea
\[
Q_d
=
U\operatorname{diag}
(0,\ldots,0,\delta_1,\ldots,\delta_{L-d})U^\top,
\qquad \delta_j>0.
\]
La coordenada usada en el paper y en el codigo es
\begin{equation}
S(\lambda;L,d)
=
1-
\frac{1}{L-d}
\sum_{j=1}^{L-d}
\frac{1}{1+\lambda\delta_j}.
\label{eq:smoothness}
\end{equation}
Entonces
\[
S(0)=0,
\qquad
S(\infty)=1,
\]
y
\begin{equation}
\frac{\partial S}{\partial\lambda}
=
\frac{1}{L-d}
\sum_{j=1}^{L-d}
\frac{\delta_j}{(1+\lambda\delta_j)^2}
>0.
\label{eq:smonotone}
\end{equation}
Por tanto, \(S\) y \(\lambda\) ordenan los modelos de la misma manera, pero
\(S\) vive en el intervalo compacto \([0,1]\). Escribimos
\[
\lambda=\lambda(S;L,d)
\]
para la inversa numerica de \eqref{eq:smoothness}.

\section{Como se continua la tendencia hacia el futuro}

El pronostico no es un polinomio ajustado por una regresion separada del
suavizador. Se toma la tendencia penalizada al final de la ventana y se
continua imponiendo
\begin{equation}
\D^d\widehat{\tau}_t=0
\qquad
\text{para los puntos futuros}.
\label{eq:zero-difference}
\end{equation}
La recurrencia usada por el codigo es
\begin{equation}
\widehat{\tau}_{n+1}
=
-
\sum_{k=0}^{d-1}
(-1)^{d-k}
\binom{d}{k}
\widehat{\tau}_{n-d+1+k}.
\label{eq:continuation}
\end{equation}
Esto produce una continuacion constante para \(d=1\), lineal para \(d=2\),
cuadratica para \(d=3\) y cubica para \(d=4\), en el sentido de diferencias
finitas. Por eso \(d\) determina tambien la geometria de la extrapolacion.

\section{Dos validaciones consecutivas}

Para un rolling origin \(r\), usamos tres bloques consecutivos:
\[
\boxed{\mathcal{T}_r}
\;\longrightarrow\;
\boxed{\mathcal{V}^{(1)}_r}
\;\longrightarrow\;
\boxed{\mathcal{V}^{(2)}_r},
\]
con
\[
|\mathcal{T}_r|=L,
\qquad
|\mathcal{V}^{(1)}_r|
=
|\mathcal{V}^{(2)}_r|
=
h.
\]

Validation 1 descubre la geometria local de candidatos. Validation 2 mide si
cada candidato sigue funcionando en el bloque temporal inmediatamente
posterior.

Para \(S\), sea
\[
\widehat{\bm z}^{(1)}_{r,d,L,h}(S)
\]
el forecast de Validation 1 y \(\bm z^{(1)}_r\) su valor observado. La perdida
es
\begin{equation}
F^{(1)}_{r,d,L,h}(S)
=
\frac{1}{h}
\left\|
\bm z^{(1)}_r
-
\widehat{\bm z}^{(1)}_{r,d,L,h}(S)
\right\|_2^2.
\label{eq:val1}
\end{equation}
La funcion \(S\mapsto F^{(1)}_{r,d,L,h}(S)\) puede tener varios minimos
locales.

\section{Paso 0: seleccionar la longitud de ventana}

Antes de seguir ramas, el codigo escoge una sola longitud \(L\) para cada
orden \(d\). Sea \(\mathcal{R}_{d,L}\) el conjunto de rolling origins historicos
para los que caben Validation 1 y Validation 2 completas. Se define
\begin{equation}
\overline F^{(1)}_{d,L}(S)
=
\frac{1}{|\mathcal{R}_{d,L}|}
\sum_{r\in\mathcal{R}_{d,L}}
F^{(1)}_{r,d,L,h}(S).
\label{eq:aggregate}
\end{equation}

Para cada \(L\), sea
\[
\widehat S_{d,L}^{\mathrm{agg}}
\in
\argmin_{S\in\mathcal{C}_{d,L}^{\mathrm{agg}}}
\overline F^{(1)}_{d,L}(S),
\]
donde \(\mathcal{C}_{d,L}^{\mathrm{agg}}\) contiene los minimos recuperados y
los endpoints relevantes. Entonces
\begin{equation}
L_d^\star
=
\argmin_L
\overline F^{(1)}_{d,L}
\left(
\widehat S_{d,L}^{\mathrm{agg}}
\right).
\label{eq:Lselection}
\end{equation}

Este paso no escoge todavia el smoothness final. Solo fija la escala temporal
sobre la cual se seguiran los minimos.

\section{Paso 1: inicializar hasta cinco minimos locales}

Fijados \(d\) y \(L_d^\star\), en el primer rolling origin se define
\begin{equation}
\mathcal{M}_{d,1}
=
\left\{
S\in(0,1):
\frac{\partial F^{(1)}_{1,d,L_d^\star,h}}{\partial S}(S)=0,
\quad
\frac{\partial^2 F^{(1)}_{1,d,L_d^\star,h}}{\partial S^2}(S)>0
\right\},
\label{eq:firstminima}
\end{equation}
mas los endpoints que actuan como minimos de frontera.

Se deduplican soluciones a menos de \(0.02\) en la escala \(S\). Si hay mas de
cinco, se conservan los cinco con menor perdida de Validation 1:
\[
S_{d,1,1},\ldots,S_{d,J_d,1},
\qquad
J_d\le5.
\]
Cada uno inicia una rama \(\mathcal{B}_{d,j}\).

\section{Paso 2: evaluar cada minimo en Validation 2}

Para cada \(S_{d,j,r}\), Validation 1 ya fue utilizada para descubrir el
candidato. Para evaluarlo predictivamente, el modelo se refitea usando todos
los datos disponibles hasta el final de Validation 1, manteniendo
\[
d,
\qquad
L_d^\star,
\qquad
S_{d,j,r}
\]
fijos. Sea \(\widehat{\bm z}^{(2)}_{d,j,r}\) el forecast del siguiente bloque.

La metrica default es RMSE en nivel:
\begin{equation}
E^{(2)}_{d,j,r}
=
\sqrt{
\frac{1}{h}
\sum_{\ell=1}^{h}
\left(
Y^{(2)}_{r,\ell}
-
\widehat Y^{(2)}_{d,j,r,\ell}
\right)^2
}.
\label{eq:val2level}
\end{equation}
Tambien se guarda RMSE en log-escala:
\begin{equation}
E^{(2,\log)}_{d,j,r}
=
\sqrt{
\frac{1}{h}
\sum_{\ell=1}^{h}
\left(
X^{(2)}_{r,\ell}
-
\widehat X^{(2)}_{d,j,r,\ell}
\right)^2
}.
\label{eq:val2log}
\end{equation}

La separacion conceptual es
\[
\boxed{F^{(1)}\text{ genera candidatos}}
\qquad
\text{y}
\qquad
\boxed{E^{(2)}\text{ juzga su utilidad futura}}.
\]

\section{Paso 3: mover la ventana y volver a buscar minimos}

Avanzamos el rolling origin una cantidad pequena \(q\). En el codigo,
\(q=1\) para GDP y \(q=5\) para las series diarias. Se vuelve a construir
\[
F^{(1)}_{r+1,d,L_d^\star,h}(S)
\]
y se recupera su conjunto de minimos locales
\[
\mathcal{M}_{d,r+1}.
\]
La hipotesis operativa es que, si la ventana cambia poco, un minimo local bien
definido deberia moverse poco en \(S\).

\section{Paso 4: seguimiento local de cada minimo}

Supongamos que la rama \(j\) tenia ultimo minimo conocido \(S_{d,j,r}\). Un
candidato \(\widetilde S\in\mathcal{M}_{d,r+1}\) puede continuar esa rama solo
si
\begin{equation}
|\widetilde S-S_{d,j,r}|
\le
\varepsilon_{\mathrm{track}}.
\label{eq:track}
\end{equation}

Entre todos los pares rama--candidato se consideran las distancias
\[
D_{jk}
=
|S_{d,j,r}-\widetilde S_{d,k,r+1}|.
\]
El codigo ordena los pares por distancia y realiza una asignacion greedy
uno-a-uno: toma primero el par mas cercano permitido, despues el siguiente que
no reutilice ni la rama ni el candidato, y asi sucesivamente.

La regla implementada aproxima
\begin{equation}
S_{d,j,r+1}
\approx
\argmin_{\widetilde S\in\mathcal{M}_{d,r+1}}
|\widetilde S-S_{d,j,r}|
\quad
\text{sujeto a}
\quad
|\widetilde S-S_{d,j,r}|
\le
\varepsilon_{\mathrm{track}},
\label{eq:continuation-rule}
\end{equation}
con correspondencia uno-a-uno entre ramas y minimos nuevos.

Si ninguna raiz nueva cae dentro del radio permitido, se registra
\[
\texttt{status}=\texttt{missing\_local\_minimum}.
\]
No se sustituye automaticamente por el mejor minimo global. La rama conserva
como referencia su ultimo \(S\) exitosamente observado y puede reaparecer mas
adelante si vuelve a existir un minimo dentro del vecindario local.

\section{Por que el seguimiento local tiene sentido}

Imaginemos provisionalmente que el origen temporal puede representarse por un
parametro continuo \(u\). Sea \(F(S,u)\) la perdida de Validation 1 cuando la
ventana esta localizada alrededor de \(u\). Supongamos que
\[
F_S(S_0,u_0)=0,
\qquad
F_{SS}(S_0,u_0)>0.
\]
Si \(F\) es suficientemente suave y la curvatura no se anula, el teorema de la
funcion implicita da una trayectoria local \(S^\star(u)\) con
\[
F_S(S^\star(u),u)=0,
\]
y
\begin{equation}
\frac{dS^\star(u)}{du}
=
-
\frac{F_{Su}(S^\star(u),u)}
{F_{SS}(S^\star(u),u)}.
\label{eq:implicit}
\end{equation}

Por tanto, si la ventana cambia poco, el termino cruzado no es excesivamente
grande y la curvatura permanece separada de cero,
\[
S^\star(u+\delta)
\approx
S^\star(u)+O(\delta).
\]
La condicion \eqref{eq:track} es una version discreta de esta idea de
continuacion de soluciones.

Tambien explica por que no conviene promediar globalmente minimos de cuencas
distintas. Si dos ramas persistentes estan cerca de \(0.10\) y \(0.90\), su
media es \(0.50\), pero nada garantiza que \(0.50\) sea un minimo. El objeto
relevante es cada trayectoria local
\[
\mathcal{B}_{d,j}
=
(S_{d,j,1},S_{d,j,2},\ldots).
\]

\section{Paso 5: almacenar la historia de cada rama}

Para cada rolling origin y rama se guarda
\[
S_{d,j,r},
\qquad
F^{(1)}_{r,d,L_d^\star,h}(S_{d,j,r}),
\qquad
E^{(2)}_{d,j,r},
\]
junto con
\[
\Delta S_{d,j,r}
=
|S_{d,j,r}-S_{d,j,r-1}|.
\]
Tambien se guarda el numero total de minimos de la superficie,
\[
M_{d,r}=|\mathcal{M}_{d,r}|,
\]
el numero de ramas emparejadas, las fechas exactas de ambos bloques y el
forecast completo sobre Validation 2.

\section{Paso 6: persistencia de una rama}

Sea \(R_{d,j}\) el numero total de rolling origins considerados para la rama
\(j\), y sea \(R^{\mathrm{match}}_{d,j}\) el numero de origins en que se encontro
una continuacion local. La fraccion de soporte es
\begin{equation}
P_{d,j}
=
\frac{R^{\mathrm{match}}_{d,j}}{R_{d,j}}.
\label{eq:support}
\end{equation}

Una rama con \(P_{d,j}=1\) persistio en todos los origins historicos. Esta
cantidad actua como regularizacion de estabilidad: una rama con un RMSE
ocasionalmente bajo pero que desaparece con frecuencia no se trata igual que
una rama que puede seguirse coherentemente.

\section{Paso 7: score historico de Validation 2}

Para cada rama se promedian unicamente los scores de los origins en los que la
rama fue emparejada. Con la metrica default,
\begin{equation}
A_{d,j}
=
\frac{1}{R^{\mathrm{match}}_{d,j}}
\sum_{r:\,(d,j)\text{ existe en }r}
E^{(2)}_{d,j,r}.
\label{eq:branchscore}
\end{equation}
Con \texttt{--selection-metric log\_rmse} se reemplaza \(E^{(2)}\) por
\(E^{(2,\log)}\).

\section{Paso 8: ultima Validation 1 antes del test}

Toda la historia usada para tracking termina antes de los ultimos \(2h\) puntos.
La estructura final es
\[
\boxed{\text{historia para tracking}}
\;\longrightarrow\;
\boxed{\mathcal{V}^{(1)}_{\mathrm{final}}}
\;\longrightarrow\;
\boxed{\mathcal{T}_{\mathrm{test}}},
\]
donde los dos ultimos bloques tienen longitud \(h\).

En \(\mathcal{V}^{(1)}_{\mathrm{final}}\) se construye una ultima superficie de
Validation 1 para cada \(d\). Cada rama historica se intenta continuar una vez
mas:
\begin{equation}
S_{d,j,\mathrm{final}}
=
\argmin_{S\in\mathcal{M}_{d,\mathrm{final}}}
|S-S_{d,j,\mathrm{last}}|
\quad
\text{sujeto a}
\quad
|S-S_{d,j,\mathrm{last}}|
\le
\varepsilon_{\mathrm{track}}.
\label{eq:final-continuation}
\end{equation}
Una rama que no puede continuar hasta esta superficie final no es elegible para
el forecast final.

\section{Paso 9: seleccion de la rama ganadora}

Sea \(\mathcal{J}_{\mathrm{final}}\) el conjunto de ramas que pudieron continuar
hasta la superficie final. El codigo aplica primero una regla de persistencia.

Si existe al menos una rama con \(P_{d,j}=1\), el conjunto elegible es
\begin{equation}
\mathcal{E}
=
\left\{
(d,j)\in\mathcal{J}_{\mathrm{final}}:
P_{d,j}=1
\right\}.
\label{eq:eligible1}
\end{equation}

Si ninguna tiene soporte completo, sea
\[
P_{\max}
=
\max_{(d,j)\in\mathcal{J}_{\mathrm{final}}}P_{d,j},
\]
y
\begin{equation}
\mathcal{E}
=
\left\{
(d,j)\in\mathcal{J}_{\mathrm{final}}:
P_{d,j}=P_{\max}
\right\}.
\label{eq:eligible2}
\end{equation}

Dentro de ese conjunto, la rama ganadora es
\begin{equation}
(d^\star,j^\star)
=
\argmin_{(d,j)\in\mathcal{E}}
A_{d,j}.
\label{eq:winner}
\end{equation}

En caso de empate numerico, la implementacion actual rompe el empate por
\(d\) y despues por el identificador de rama. La regla puede resumirse como
\[
\boxed{
\text{persistencia primero}
\quad+\quad
\text{menor perdida historica en Validation 2}
}.
\]

\section{Cual es el smoothness finalmente seleccionado}

Una vez seleccionada la rama \((d^\star,j^\star)\), el smoothness utilizado para
el forecast final es
\begin{equation}
\boxed{
S_{\mathrm{final}}^\star
=
S_{d^\star,j^\star,\mathrm{final}}.
}
\label{eq:finalS}
\end{equation}

Es decir:
\begin{itemize}[leftmargin=2em]
    \item no es el \(S\) del primer rolling origin;
    \item no es el minimo global de todas las superficies historicas;
    \item no es el promedio de todos los \(S\);
    \item no usa el test verdadero;
    \item es el minimo local actual de la rama historicamente seleccionada.
\end{itemize}

La rama se escoge por su comportamiento historico en Validation 2, pero su
posicion final se actualiza localmente con la informacion mas reciente
disponible antes del test.

\section{Paso 10: refit final y test intacto}

Sea \(L^\star=L_{d^\star}^\star\). Tomamos las ultimas \(L^\star\)
observaciones disponibles antes del test, incluyendo el bloque final de
Validation 1, y ajustamos
\begin{equation}
\widehat{\bm{\tau}}_{\mathrm{final}}
=
\argmin_{\bm{\tau}}
\left\{
\|\bm{x}_{\mathrm{final}}-\bm{\tau}\|_2^2
+
\lambda(S_{\mathrm{final}}^\star)
\|D_{d^\star}\bm{\tau}\|_2^2
\right\}.
\label{eq:finalfit}
\end{equation}
Despues se continua esa tendencia \(h\) pasos usando
\[
\D^{d^\star}\widehat{\tau}_t=0.
\]
Solo entonces se compara con el test real.

El test no participa en la seleccion de \(L\), deteccion de minimos, tracking,
score historico \(A_{d,j}\), eleccion de rama ni seleccion de
\(S_{\mathrm{final}}^\star\). Solo produce metricas retrospectivas de
diagnostico.

\section{Por que Validation 1 y Validation 2 deben estar separadas}

Si el mismo bloque que define \(S_{d,j,r}\) tambien se utilizara para decidir
cual minimo es predictivamente mejor, la seleccion reutilizaria la misma
realizacion que genero los candidatos.

Aqui,
\[
\mathcal{V}^{(1)}:
\text{descubrimiento geometrico},
\qquad
\mathcal{V}^{(2)}:
\text{evaluacion temporal posterior}.
\]
Puede ocurrir que
\[
F^{(1)}(S_a)<F^{(1)}(S_b)
\]
pero
\[
E^{(2)}(S_a)>E^{(2)}(S_b).
\]
Ese desacuerdo es justamente informacion util: dos minimos de la misma
superficie pueden tener capacidad de continuacion muy diferente.

\section{Interpretacion jerarquica}

El metodo puede verse como
\[
\boxed{
d
\to
L_d^\star
\to
\text{ramas locales de }S
\to
\text{persistencia}
\to
\text{Validation 2}
\to
S_{\mathrm{final}}^\star
}.
\]

Cada etapa controla un aspecto distinto:
\begin{itemize}[leftmargin=2em]
    \item \(d\) determina la familia geometrica de extrapolacion;
    \item \(L\) determina la escala temporal de entrenamiento;
    \item los minimos locales de \(S\) representan escalas de suavidad
    competitivas;
    \item el tracking conserva la identidad de esas escalas a traves del tiempo;
    \item la persistencia penaliza soluciones temporalmente fragiles;
    \item Validation 2 mide utilidad predictiva posterior;
    \item la ultima Validation 1 actualiza la posicion de la rama antes del test.
\end{itemize}

Una rama no es una distribucion posterior ni un intervalo de confianza. Es una
trayectoria determinista de minimos locales de problemas de optimizacion
cercanos.

\section{Que significa un promedio local de smoothness}

El codigo actual no usa un promedio de \(S\) para el forecast final: usa
\eqref{eq:finalS}. Sin embargo, una extension coherente seria promediar
\emph{solo dentro de la misma rama y localmente en el tiempo}. Por ejemplo,
\begin{equation}
\widetilde S_{d,j,R}^{(K)}
=
\frac{1}{K}
\sum_{k=0}^{K-1}
S_{d,j,R-k},
\label{eq:localmean}
\end{equation}
o una media ponderada
\begin{equation}
\widetilde S_{d,j,R}^{(\rho)}
=
\frac{
\sum_{k=0}^{K-1}
\rho^k S_{d,j,R-k}
}{
\sum_{k=0}^{K-1}\rho^k
},
\qquad 0<\rho\le1.
\label{eq:weightedmean}
\end{equation}

Esto es muy distinto de promediar minimos de cuencas diferentes. En
\eqref{eq:localmean}, todos los terminos ya fueron identificados como la misma
rama mediante \eqref{eq:track}. El codigo reporta medias y medianas globales y
de los cinco matches mas recientes por rama para diagnostico, pero el default
actual usa el minimo local mas reciente.

\section{Limitaciones matematicas del tracking actual}

\subsection{Matching greedy}

La asignacion rama--minimo es greedy por distancia. Si dos ramas se acercan o
cruzan, la correspondencia puede depender del orden de las distancias. Una
alternativa mas formal seria resolver
\begin{equation}
\min_{\pi}
\sum_j
|S_{d,j,r}-\widetilde S_{d,\pi(j),r+1}|
\label{eq:assignment}
\end{equation}
sujeto a
\[
|S_{d,j,r}-\widetilde S_{d,\pi(j),r+1}|
\le
\varepsilon_{\mathrm{track}},
\]
por ejemplo con Hungarian matching. El codigo actual usa la aproximacion
greedy.

\subsection{Eleccion de \(\varepsilon_{\mathrm{track}}\)}

El radio \(0.10\) es una regla operativa, no un parametro inferido
estadisticamente. Si es demasiado pequeno, una rama puede fragmentarse; si es
demasiado grande, dos cuencas pueden confundirse. Por ello conviene inspeccionar
\[
\max_r|S_{d,j,r+1}-S_{d,j,r}|,
\]
la distribucion de \(\Delta S\), el numero de minimos por superficie y el
numero de matches por origin.

\subsection{Numero variable de minimos}

No se impone que
\[
|\mathcal{M}_{d,r}|
=
|\mathcal{M}_{d,r+1}|.
\]
Que el numero permanezca estable es una hipotesis empirica que puede examinarse
en \texttt{rolling\_minimum\_tracks.csv} y en el dashboard. El nacimiento,
desaparicion o fusion aparente de minimos puede ser informacion relevante.

\subsection{Seleccion de \(L\) separada del tracking}

\(L_d^\star\) se selecciona una vez por orden mediante la perdida agregada de
Validation 1. No se sigue una trayectoria temporal de \(L\). Una extension
mayor podria seguir simultaneamente ramas en \((L,S)\), pero seria un problema
de correspondencia bidimensional distinto.

\section{Criterio completo en una sola expresion}

Aunque el procedimiento es secuencial, puede resumirse como
\begin{equation}
(d^\star,j^\star)
=
\argmin_{(d,j)}
\left\{
A_{d,j}:
(d,j)\text{ es suficientemente persistente y continua al origin final}
\right\},
\label{eq:concept}
\end{equation}
donde
\[
A_{d,j}
=
\frac{1}{R^{\mathrm{match}}_{d,j}}
\sum_r E^{(2)}_{d,j,r},
\]
y la secuencia \(\{S_{d,j,r}\}_r\) satisface
\[
S_{d,j,r}\in\mathcal{M}_{d,r},
\qquad
|S_{d,j,r}-S_{d,j,r-1}|
\le
\varepsilon_{\mathrm{track}}.
\]
Despues,
\[
S_{\mathrm{final}}^\star
=
\text{continuacion local final de la rama }(d^\star,j^\star).
\]

Por eso el metodo no selecciona simplemente un numero \(S\): selecciona una
\textbf{rama estable de soluciones locales} y usa su estado mas reciente.

\section{Correspondencia con los archivos del experimento}

\begin{center}
\begin{tabular}{p{0.37\textwidth} p{0.54\textwidth}}
\toprule
Archivo & Objeto matematico principal\\
\midrule
\texttt{order\_window\_selection.csv}
&
\(L_d^\star\) y la perdida agregada de Validation 1
\\
\texttt{rolling\_minimum\_tracks.csv}
&
\(S_{d,j,r}\), \(F^{(1)}_{d,j,r}\), \(E^{(2)}_{d,j,r}\),
\(\Delta S_{d,j,r}\), soporte y conteo de minimos
\\
\texttt{rolling\_validation\_paths.csv}
&
\(\widehat{\bm z}^{(2)}_{d,j,r}\) completo para cada rama y origin
\\
\texttt{tracked\_branch\_selection.csv}
&
\(P_{d,j}\), \(A_{d,j}\), resumenes de \(S\), continuidad final y rama ganadora
\\
\texttt{tracked\_branch\_test\_paths.csv}
&
pronostico final de cada rama que pudo continuar
\\
\texttt{run\_metadata.json}
&
protocolo, parametros, hashes y rol del test
\\
\bottomrule
\end{tabular}
\end{center}

\section{Resumen del algoritmo}

Para cada serie:
\begin{enumerate}[leftmargin=2.4em]
    \item Reservar el test final de longitud \(h\).
    \item Para cada \(d\), escoger \(L_d^\star\) por
    \eqref{eq:Lselection}.
    \item En el primer rolling origin, encontrar hasta cinco minimos locales de
    Validation 1.
    \item Inicializar una rama por minimo.
    \item Evaluar cada rama en Validation 2.
    \item Mover el rolling origin y reconstruir la superficie de Validation 1.
    \item Continuar cada rama al minimo nuevo mas cercano dentro de
    \(\varepsilon_{\mathrm{track}}\), sin reutilizar minimos.
    \item Repetir y acumular \(E^{(2)}_{d,j,r}\).
    \item Calcular soporte \(P_{d,j}\) y score historico \(A_{d,j}\).
    \item Construir una ultima Validation 1 inmediatamente antes del test.
    \item Continuar localmente cada rama una vez mas.
    \item Elegir la rama mas persistente y, entre las elegibles, la de menor
    \(A_{d,j}\).
    \item Usar su nuevo minimo local como \(S_{\mathrm{final}}^\star\).
    \item Refit final y forecast del test verdadero.
\end{enumerate}

En simbolos:
\[
F^{(1)}_r(S)
\to
\mathcal{M}_r
\to
S_{j,r}
\to
E^{(2)}_{j,r}
\to
S_{j,r+1}
\to
E^{(2)}_{j,r+1}
\to
\cdots
\]
y finalmente
\[
\{E^{(2)}_{j,r}\}_r
\to
j^\star
\to
S_{\mathrm{final}}^\star
\to
\text{true test}.
\]

\section{Conclusion}

El procedimiento separa tres ideas que conviene no mezclar:
\[
\text{geometria de la perdida}
\quad
\longrightarrow
\quad
\text{persistencia de un minimo local}
\quad
\longrightarrow
\quad
\text{calidad predictiva posterior}.
\]

Validation 1 crea los candidatos; el rolling temporal los convierte en ramas;
Validation 2 asigna a cada rama una historia de desempeno fuera de la ventana
que genero el minimo. La rama ganadora debe ser suficientemente persistente y
tener el menor score historico de Validation 2 entre las ramas elegibles. El
smoothness final es su minimo local mas reciente:
\[
\boxed{
S_{\mathrm{final}}^\star
=
S_{d^\star,j^\star,\mathrm{final}}.
}
\]

Esta construccion tiene sentido porque conserva la identidad de cada cuenca de
optimizacion, evita promediar soluciones que pertenecen a minimos diferentes y
separa el descubrimiento de candidatos de su evaluacion predictiva. El test
verdadero permanece fuera de todo el proceso de seleccion.

\end{document}
